\documentclass[11pt,a4paper]{article}

% ============================================================
%  QIMMA -- COURS : Calcul des probabilités
%  2ème Bac Sciences Mathématiques (A et B)
%  Standard exhaustif : définitions formelles + démonstrations
%  Basé sur templates/qimma-template.tex (charte Qimma)
% ============================================================

\usepackage[french]{babel}
\usepackage[utf8]{inputenc}
\usepackage[T1]{fontenc}
\usepackage{lmodern}
\usepackage[a4paper,margin=2cm,top=2.3cm,bottom=2.6cm]{geometry}
\usepackage{amsmath,amssymb}
\usepackage{xcolor}
\usepackage{graphicx}
\usepackage{tikz}
\usetikzlibrary{shadows,calc}
\usepackage[most]{tcolorbox}
\usepackage{titlesec}
\usepackage{fancyhdr}
\usepackage{enumitem}
\usepackage{pifont}
\usepackage{array}
\usepackage{colortbl}
\usepackage{hyperref}

% ------------------- CHARTE QIMMA -------------------
\definecolor{qteal}{HTML}{0d9488}
\definecolor{qtealdark}{HTML}{134e4a}
\definecolor{qamber}{HTML}{fbbf24}
\definecolor{qambertext}{HTML}{92400e}
\definecolor{ink}{HTML}{1f2937}
\definecolor{inkgray}{HTML}{6b7280}
\definecolor{tealpale}{HTML}{e6f5f3}
\colorlet{colDef}{qteal}
\definecolor{colProp}{HTML}{0f766e}
\definecolor{colEx}{HTML}{b45309}
\definecolor{colRem}{HTML}{9d174d}
\color{ink}

\graphicspath{{../../../../assets/}}
\newcommand{\logofile}{../../../../assets/qimma-logo.pdf}

% ------------------- METADONNEES -------------------
\newcommand{\matiere}{Mathématiques}
\newcommand{\niveau}{2\up{ème} Bac -- Sciences Mathématiques}
\newcommand{\chapitretitre}{Calcul des probabilités}

% ------------------- EN-TETE / PIED DE PAGE -------------------
\pagestyle{fancy}
\fancyhf{}
\renewcommand{\headrulewidth}{0pt}
\renewcommand{\footrulewidth}{0.5pt}
\fancyfoot[L]{\raisebox{-0.15cm}{\includegraphics[height=0.35cm]{\logofile}}~\small\sffamily\color{inkgray}Qimma}
\fancyfoot[C]{\small\sffamily\color{inkgray}\matiere\ -- \niveau}
\fancyfoot[R]{\small\sffamily\color{qtealdark}\thepage}

% ------------------- TITRES DE SECTION -------------------
\titleformat{\section}
  {\normalfont\sffamily\Large\bfseries\color{qtealdark}}
  {\tikz[baseline=-3.2pt]{\node[circle,fill=qteal,text=white,inner sep=0pt,minimum size=8mm]{\sffamily\bfseries\thesection};}}
  {10pt}{}
  [\vspace{-2mm}{\color{qamber}\titlerule[1.6pt]}\vspace{2mm}]
\titlespacing*{\section}{0pt}{16pt}{10pt}

\titleformat{\subsection}
  {\normalfont\sffamily\large\bfseries\color{qtealdark}}
  {\color{qteal}\ding{212}}{6pt}{}
\titlespacing*{\subsection}{0pt}{12pt}{6pt}

% ------------------- BOITES PEDAGOGIQUES -------------------
\newtcolorbox{fiche}[3][]{
  enhanced, breakable,
  colback=white, colframe=#2,
  boxrule=1pt, arc=2mm,
  top=7mm, left=4mm, right=4mm, bottom=3mm,
  drop shadow={black!25, shadow xshift=1mm, shadow yshift=-1mm},
  attach boxed title to top left={xshift=4mm, yshift=-3mm, yshifttext=-1mm},
  boxed title style={colback=#2, arc=1.5mm, boxrule=0pt, left=2mm, right=2mm},
  coltitle=white, fonttitle=\sffamily\bfseries,
  title={#3}, #1
}
\newenvironment{definition}[1][Définition]
  {\begin{fiche}{colDef}{\ding{110}~~#1}}{\end{fiche}}
\newenvironment{propriete}[1][Propriété]
  {\begin{fiche}{colProp}{\ding{72}~~#1}}{\end{fiche}}
\newenvironment{exemple}[1][Exemple]
  {\begin{fiche}{colEx}{\ding{217}~~#1}}{\end{fiche}}
\newenvironment{remarque}[1][Remarque]
  {\begin{fiche}{colRem}{\ding{43}~~#1}}{\end{fiche}}

% Démonstration (hors boîte, style discret)
\newenvironment{demo}
  {\par\smallskip\noindent{\sffamily\bfseries\color{qtealdark}Démonstration.}\ \itshape}
  {\ \normalfont\hfill$\blacksquare$\par\medskip}

\hypersetup{colorlinks=true, linkcolor=qtealdark, urlcolor=qtealdark}

% ============================================================
\begin{document}

% ------------------- BANNIERE DE TITRE -------------------
\begin{tcolorbox}[
  enhanced, boxrule=0pt, arc=4mm,
  interior style={left color=qtealdark, right color=qteal},
  top=8mm, bottom=8mm, left=10mm, right=32mm,
  drop shadow={black!35, shadow xshift=1.5mm, shadow yshift=-1.5mm},
  before skip=0pt, after skip=9mm,
  overlay={
    \node[anchor=north west] at ([xshift=6mm,yshift=-6mm]frame.north west)
      {\includegraphics[height=1.25cm]{\logofile}};
    \node[anchor=north east, fill=qamber, text=qambertext, rounded corners=2pt,
          inner sep=4pt, font=\sffamily\bfseries\small] at ([xshift=-6mm,yshift=-6mm]frame.north east) {COURS};
  }
]
\hspace{1.6cm}\centering
{\sffamily\bfseries\color{white}\fontsize{23}{27}\selectfont \chapitretitre}\\[3mm]
{\sffamily\color{white!90}\large \matiere\ \textbullet\ \niveau}
\end{tcolorbox}

% ------------------- OBJECTIFS -------------------
\begin{tcolorbox}[
  enhanced, colback=tealpale, colframe=qteal, boxrule=0.8pt, arc=2mm,
  left=5mm, right=5mm, top=3mm, bottom=3mm,
  title={\sffamily\bfseries\color{qtealdark}\ding{51}~~Objectifs du chapitre},
  coltitle=qtealdark, colbacktitle=tealpale, boxrule=0.8pt,
  attach title to upper={\par\vspace{1mm}}
]
\begin{itemize}[leftmargin=6mm, label=\ding{212}, itemsep=1mm]
  \item Construire le vocabulaire probabiliste (univers, événements) et démontrer les axiomes et propriétés de base.
  \item Calculer une probabilité dans le cas de l'équiprobabilité, en combinant avec les outils de dénombrement.
  \item Maîtriser les probabilités conditionnelles, l'indépendance d'événements, et la formule des probabilités totales.
  \item Reconnaître et utiliser une variable aléatoire, sa loi de probabilité, son espérance, sa variance et son écart-type.
  \item Déterminer et représenter la fonction de répartition d'une variable aléatoire discrète.
  \item Connaître la loi binomiale : construction, formule, espérance et variance.
  \item Traiter les exercices type examen national : arbres pondérés, urnes, lois usuelles, prise de décision.
\end{itemize}
\end{tcolorbox}

\vspace{4mm}

% ------------------- SECTION 1 -------------------
\section{Vocabulaire et axiomes des probabilités}

\begin{definition}[Univers, événements]
Une expérience aléatoire a pour \textbf{univers} $\Omega$ l'ensemble de ses résultats (issues) possibles. Un \textbf{événement} est une partie de $\Omega$. Un événement à un seul élément est \textbf{élémentaire}. $\varnothing$ est l'\textbf{événement impossible}, $\Omega$ l'\textbf{événement certain}.
\end{definition}

\begin{definition}[Probabilité]
Une \textbf{probabilité} sur $\Omega$ fini est une application $P$ qui à chaque événement $A\subset\Omega$ associe $P(A)\in[0;1]$, vérifiant :
\[
  P(\Omega)=1, \qquad P(A\cup B)=P(A)+P(B) \quad\text{si } A\cap B=\varnothing.
\]
\end{definition}

\begin{propriete}[Conséquences des axiomes]
\[
  P(\varnothing)=0, \qquad P(\bar A)=1-P(A), \qquad P(A\cup B)=P(A)+P(B)-P(A\cap B),
\]
\[
  A\subset B \implies P(A)\leq P(B).
\]
\end{propriete}

\begin{demo}
$P(A\cup B)=P(A)+P(B)-P(A\cap B)$ : on décompose $A\cup B=A\cup(B\setminus A)$, union disjointe, donc $P(A\cup B)=P(A)+P(B\setminus A)$. Or $B=(B\cap A)\cup(B\setminus A)$, union disjointe, donc $P(B)=P(A\cap B)+P(B\setminus A)$, soit $P(B\setminus A)=P(B)-P(A\cap B)$. En substituant : $P(A\cup B)=P(A)+P(B)-P(A\cap B)$.
\end{demo}

\begin{definition}[Équiprobabilité]
Si $\Omega$ est fini et que tous les résultats élémentaires sont \textbf{équiprobables} (même probabilité $\frac{1}{\text{card}(\Omega)}$), alors pour tout événement $A$ :
\[
  P(A) = \frac{\text{card}(A)}{\text{card}(\Omega)} = \frac{\text{nombre de cas favorables}}{\text{nombre de cas possibles}}.
\]
\end{definition}

\begin{exemple}[Calcul en situation d'équiprobabilité]
On tire une carte au hasard dans un jeu de $32$ cartes. Calculons $P(\text{« tirer un roi »})$ et $P(\text{« tirer un coeur »})$.
\[
  P(\text{roi}) = \frac{4}{32} = \frac18, \qquad P(\text{coeur}) = \frac{8}{32} = \frac14.
\]
\end{exemple}

% ------------------- SECTION 2 -------------------
\section{Probabilités conditionnelles}

\begin{definition}[Probabilité conditionnelle]
Soient $A,B$ événements avec $P(A)>0$. La \textbf{probabilité de $B$ sachant $A$} est
\[
  P_A(B) = \frac{P(A\cap B)}{P(A)}.
\]
\end{definition}

\begin{propriete}[Formule des probabilités composées]
Si $P(A)>0$ : $P(A\cap B)=P(A)\times P_A(B)$.
\end{propriete}

\begin{definition}[Événements indépendants]
$A$ et $B$ sont \textbf{indépendants} lorsque $P(A\cap B)=P(A)\times P(B)$ (de façon équivalente, si $P(A)>0$ : $P_A(B)=P(B)$, la réalisation de $A$ n'influence pas celle de $B$).
\end{definition}

\begin{propriete}[Système complet d'événements et probabilités totales]
$(A_1,\ldots,A_n)$ forme un \textbf{système complet d'événements} si les $A_i$ sont deux à deux disjoints et $A_1\cup\cdots\cup A_n=\Omega$ (avec $P(A_i)>0$). Alors, pour tout événement $B$ :
\[
  P(B) = \sum_{i=1}^n P(A_i)\times P_{A_i}(B) \quad\text{(formule des probabilités totales)}.
\]
\end{propriete}

\begin{demo}
$B=(B\cap A_1)\cup\cdots\cup(B\cap A_n)$, union \textbf{disjointe} (car les $A_i$ le sont). Donc $P(B)=\sum_iP(B\cap A_i)=\sum_iP(A_i)P_{A_i}(B)$ (formule des probabilités composées).
\end{demo}

\begin{exemple}[Arbre pondéré et probabilités totales]
Une urne $U_1$ contient $3$ boules rouges et $2$ boules vertes ; une urne $U_2$ contient $1$ boule rouge et $4$ boules vertes. On choisit une urne au hasard (équiprobabilité), puis on y tire une boule. Calculons $P(\text{« tirer rouge »})$.

\smallskip
Soit $A_1$ : « choisir $U_1$ », $A_2$ : « choisir $U_2$ », $R$ : « tirer rouge ». $P(A_1)=P(A_2)=\frac12$, $P_{A_1}(R)=\frac35$, $P_{A_2}(R)=\frac15$.
\[
  P(R) = \frac12\times\frac35+\frac12\times\frac15 = \frac{3}{10}+\frac1{10} = \frac4{10} = \frac25.
\]
\end{exemple}

\begin{propriete}[Formule de Bayes]
Si $(A_1,\ldots,A_n)$ système complet, $P(B)>0$ :
\[
  P_B(A_i) = \frac{P(A_i)\times P_{A_i}(B)}{P(B)} = \frac{P(A_i)P_{A_i}(B)}{\sum_jP(A_j)P_{A_j}(B)}.
\]
\end{propriete}

\begin{exemple}[Application de Bayes -- reprise de l'exemple précédent]
Sachant que l'on a tiré une boule rouge, calculons la probabilité que ce soit l'urne $U_1$ qui ait été choisie.
\[
  P_R(A_1) = \frac{P(A_1)P_{A_1}(R)}{P(R)} = \frac{\frac12\times\frac35}{\frac25} = \frac{\frac3{10}}{\frac25} = \frac3{10}\times\frac52 = \frac34.
\]
\end{exemple}

% ------------------- SECTION 3 -------------------
\section{Variables aléatoires}

\begin{definition}[Variable aléatoire, loi de probabilité]
Une \textbf{variable aléatoire} $X$ sur $\Omega$ (fini) est une fonction $X:\Omega\to\mathbb{R}$. Si $X$ prend les valeurs $x_1,\ldots,x_n$, sa \textbf{loi de probabilité} est la donnée des $P(X=x_i)$ pour chaque $i$, avec $\sum_iP(X=x_i)=1$.
\end{definition}

\begin{definition}[Espérance, variance, écart-type]
\[
  E(X) = \sum_{i=1}^nx_iP(X=x_i), \qquad
  V(X) = \sum_{i=1}^n\bigl(x_i-E(X)\bigr)^2P(X=x_i), \qquad
  \sigma(X) = \sqrt{V(X)}.
\]
\end{definition}

\begin{propriete}[Formule de König-Huygens]
\[
  V(X) = E(X^2)-\bigl(E(X)\bigr)^2.
\]
\end{propriete}

\begin{demo}
En développant : $V(X)=\sum_i(x_i-E(X))^2P(X=x_i) = \sum_i\bigl[x_i^2-2x_iE(X)+E(X)^2\bigr]P(X=x_i)$
\[
  = \sum_ix_i^2P(X=x_i) - 2E(X)\sum_ix_iP(X=x_i) + E(X)^2\sum_iP(X=x_i).
\]
Or $\sum_ix_i^2P(X=x_i)=E(X^2)$, $\sum_ix_iP(X=x_i)=E(X)$ et $\sum_iP(X=x_i)=1$, donc
\[
  V(X) = E(X^2)-2E(X)^2+E(X)^2 = E(X^2)-E(X)^2.
\]
\end{demo}

\begin{propriete}[Linéarité de l'espérance]
Pour $a,b\in\mathbb{R}$ : $E(aX+b)=aE(X)+b$, et $V(aX+b)=a^2V(X)$.
\end{propriete}

\begin{exemple}[Loi de probabilité complète]
On lance deux dés équilibrés et on note $X$ la somme des points obtenus. Établissons la loi de $X$ et calculons $E(X)$.

\smallskip
$X$ prend les valeurs de $2$ à $12$. Par exemple $P(X=2)=\frac1{36}$ (seul $(1,1)$), $P(X=7)=\frac6{36}$ (6 combinaisons).

\begin{center}
\renewcommand{\arraystretch}{1.3}
\small
\begin{tabular}{c|ccccccccccc}
$x_i$ & 2 & 3 & 4 & 5 & 6 & 7 & 8 & 9 & 10 & 11 & 12 \\ \hline
$P(X=x_i)$ & $\frac1{36}$ & $\frac2{36}$ & $\frac3{36}$ & $\frac4{36}$ & $\frac5{36}$ & $\frac6{36}$ & $\frac5{36}$ & $\frac4{36}$ & $\frac3{36}$ & $\frac2{36}$ & $\frac1{36}$ \\
\end{tabular}
\end{center}

\[
  E(X) = \frac{1}{36}\Bigl(2\times1+3\times2+\cdots+12\times1\Bigr) = \frac{252}{36} = 7.
\]
\end{exemple}

\begin{definition}[Fonction de répartition]
Soit $X$ une variable aléatoire prenant les valeurs $x_1<x_2<\cdots<x_n$. La \textbf{fonction de répartition} de $X$ est la fonction $F_X:\mathbb{R}\to[0\,;1]$ définie par
\[
  F_X(t) = P(X\leq t) = \sum_{x_i\leq t}P(X=x_i).
\]
\end{definition}

\begin{propriete}[Propriétés de $F_X$]
\begin{itemize}[leftmargin=6mm, itemsep=0.5mm]
  \item $F_X$ est \textbf{en escalier}, \textbf{croissante} (au sens large) sur $\mathbb{R}$, constante entre deux valeurs consécutives prises par $X$.
  \item $F_X(t)=0$ pour $t<x_1$ ; \quad $F_X(t)=1$ pour $t\geq x_n$.
  \item Pour tout $i$ : $F_X(x_i)-F_X(x_{i-1}) = P(X=x_i)$ (le « saut » de $F_X$ en $x_i$ vaut $P(X=x_i)$), avec la convention $F_X(x_0)=0$.
  \item Pour $a<b$ : $P(a<X\leq b)=F_X(b)-F_X(a)$.
\end{itemize}
\end{propriete}

\begin{exemple}[Construire et représenter $F_X$]
Soit $X$ la variable aléatoire de loi : $P(X=1)=0{,}2$, $P(X=2)=0{,}5$, $P(X=3)=0{,}3$. Déterminons $F_X$.

\smallskip
\[
  F_X(t) =
  \begin{cases}
    0 & \text{si } t<1,\\
    0{,}2 & \text{si } 1\leq t<2,\\
    0{,}7 & \text{si } 2\leq t<3,\\
    1 & \text{si } t\geq3.
  \end{cases}
\]
La courbe de $F_X$ est une fonction \textbf{en escalier}, croissante par sauts de hauteur $0{,}2$ en $t=1$, $0{,}5$ en $t=2$, $0{,}3$ en $t=3$ (sauts qui totalisent bien $1$).

\smallskip
Par exemple, $P(1<X\leq3)=F_X(3)-F_X(1)=1-0{,}2=0{,}8$ (cohérent avec $P(X=2)+P(X=3)=0{,}5+0{,}3=0{,}8$).
\end{exemple}

% ------------------- SECTION 4 -------------------
\section{Loi binomiale}

\begin{definition}[Épreuve de Bernoulli, schéma de Bernoulli]
Une \textbf{épreuve de Bernoulli} de paramètre $p$ est une expérience à deux issues, « succès » (probabilité $p$) et « échec » (probabilité $1-p$). Un \textbf{schéma de Bernoulli} est la répétition de $n$ épreuves de Bernoulli \textbf{identiques et indépendantes}.
\end{definition}

\begin{propriete}[Loi binomiale]
Dans un schéma de Bernoulli de paramètres $n$ et $p$, la variable aléatoire $X$ comptant le \textbf{nombre de succès} suit la \textbf{loi binomiale} $\mathcal{B}(n,p)$ :
\[
  P(X=k) = \binom nk p^k(1-p)^{n-k}, \qquad k=0,1,\ldots,n.
\]
\end{propriete}

\begin{demo}
Une issue donnée avec exactement $k$ succès (parmi $n$ épreuves) a pour probabilité $p^k(1-p)^{n-k}$ (indépendance des épreuves). Le nombre de façons de placer ces $k$ succès parmi les $n$ épreuves est $\dbinom nk$ (choix des positions). Par incompatibilité de ces différentes issues (union disjointe), on somme : $P(X=k)=\dbinom nkp^k(1-p)^{n-k}$.
\end{demo}

\begin{propriete}[Espérance et variance de la loi binomiale]
Si $X\sim\mathcal{B}(n,p)$ :
\[
  E(X) = np, \qquad V(X) = np(1-p), \qquad \sigma(X) = \sqrt{np(1-p)}.
\]
\end{propriete}

\begin{exemple}[Application de la loi binomiale]
On lance $5$ fois un dé équilibré. Soit $X$ le nombre de fois où l'on obtient un $6$. $X\sim\mathcal{B}\left(5,\frac16\right)$.
\[
  P(X=2) = \binom52\left(\frac16\right)^2\left(\frac56\right)^3 = 10\times\frac1{36}\times\frac{125}{216} = \frac{1250}{7776} \approx 0{,}161.
\]
\[
  E(X) = 5\times\frac16 = \frac56, \qquad V(X) = 5\times\frac16\times\frac56 = \frac{25}{36}.
\]
\end{exemple}

% ------------------- SECTION 5 -------------------
\section{Exercices corrigés -- niveau examen national SM}

\begin{exemple}[Exercice 1 -- équiprobabilité et dénombrement combiné]
Une urne contient $10$ boules : $6$ rouges et $4$ noires. On tire simultanément $3$ boules. Calculer la probabilité d'obtenir exactement $2$ rouges.

\smallskip
\textbf{Corrigé.} $\text{card}(\Omega)=\dbinom{10}3=120$. Cas favorables : $2$ rouges parmi $6$ et $1$ noire parmi $4$ : $\dbinom62\times\dbinom41=15\times4=60$.
\[
  P = \frac{60}{120} = \frac12.
\]
\end{exemple}

\begin{exemple}[Exercice 2 -- indépendance à vérifier]
On lance un dé équilibré. Soient $A$ : « obtenir un nombre pair », $B$ : « obtenir un multiple de $3$ ». $A$ et $B$ sont-ils indépendants ?

\smallskip
\textbf{Corrigé.} $P(A)=\frac36=\frac12$ ($\{2,4,6\}$), $P(B)=\frac26=\frac13$ ($\{3,6\}$). $A\cap B=\{6\}$, donc $P(A\cap B)=\frac16$.

\smallskip
$P(A)\times P(B)=\frac12\times\frac13=\frac16=P(A\cap B)$. \textbf{Conclusion} : $A$ et $B$ sont \textbf{indépendants}.
\end{exemple}

\begin{exemple}[Exercice 3 -- variable aléatoire avec gain/perte]
Un jeu consiste à tirer une carte parmi $32$. On gagne $10$ DH si c'est un roi, on perd $2$ DH sinon. Soit $X$ le gain algébrique. Déterminer la loi de $X$, calculer $E(X)$ et interpréter.

\smallskip
\textbf{Corrigé.} $P(X=10)=\dfrac{4}{32}=\dfrac18$ ; $P(X=-2)=\dfrac{28}{32}=\dfrac78$.
\[
  E(X) = 10\times\frac18+(-2)\times\frac78 = \frac{10}8-\frac{14}8 = -\frac48 = -0{,}5.
\]
\textbf{Interprétation} : le jeu est défavorable au joueur (en moyenne, il perd $0{,}5$ DH par partie).
\end{exemple}

\begin{exemple}[Exercice 4 -- probabilités totales avec arbre à trois branches]
Une usine produit des pièces sur trois machines $M_1$, $M_2$, $M_3$, représentant respectivement $30\%$, $50\%$, $20\%$ de la production. Les taux de défaut sont respectivement $2\%$, $1\%$, $3\%$. On prend une pièce au hasard. Calculer la probabilité qu'elle soit défectueuse.

\smallskip
\textbf{Corrigé.} Soit $D$ : « la pièce est défectueuse ».
\[
  P(D) = 0{,}3\times0{,}02+0{,}5\times0{,}01+0{,}2\times0{,}03 = 0{,}006+0{,}005+0{,}006 = 0{,}017.
\]
\textbf{Conclusion} : $P(D)=1{,}7\%$.
\end{exemple}

\begin{exemple}[Exercice 5 -- loi binomiale, calcul « au moins »]
On répète $4$ fois, de façon indépendante, une expérience où la probabilité de succès est $p=0{,}3$. Calculer la probabilité d'obtenir au moins un succès.

\smallskip
\textbf{Corrigé.} Soit $X\sim\mathcal{B}(4;0{,}3)$. Par le complémentaire :
\[
  P(X\geq1) = 1-P(X=0) = 1-\binom40(0{,}3)^0(0{,}7)^4 = 1-0{,}7^4 = 1-0{,}2401 = 0{,}7599.
\]
\end{exemple}

% ------------------- SECTION 6 -------------------
\section{Méthodes et pièges : la boîte à outils de l'examen}

\subsection{Ce qu'on te demande, ce que tu fais}

\begin{center}
\renewcommand{\arraystretch}{1.45}
\sffamily\small
\begin{tabular}{>{\raggedright\arraybackslash}p{7.2cm} >{\raggedright\arraybackslash}p{8cm}}
\rowcolor{qtealdark} \color{white}\bfseries Ce qu'on te demande & \color{white}\bfseries Ton réflexe \\
\rowcolor{tealpale} Calculer $P(A)$ en équiprobabilité & $\frac{\text{card}(A)}{\text{card}(\Omega)}$, avec dénombrement (arrangement/combinaison) pour les cardinaux. \\
Situation en deux étapes (urne puis tirage...) & Arbre pondéré : produit sur chaque branche, somme sur les branches menant au résultat. \\
\rowcolor{tealpale} « Sachant que... » & Probabilité conditionnelle $P_A(B)=\frac{P(A\cap B)}{P(A)}$. \\
Retrouver la cause à partir d'un effet observé & Formule de Bayes. \\
\rowcolor{tealpale} Vérifier une indépendance & Comparer $P(A\cap B)$ à $P(A)\times P(B)$. \\
Loi d'une variable aléatoire & Lister toutes les valeurs possibles, calculer chaque probabilité, vérifier que la somme vaut $1$. \\
\rowcolor{tealpale} Déterminer $F_X$ & Fonction en escalier : $F_X(t)=\sum_{x_i\leq t}P(X=x_i)$, sauts de hauteur $P(X=x_i)$ en chaque $x_i$. \\
Répétitions identiques indépendantes & Reconnaître un schéma de Bernoulli, utiliser $\mathcal{B}(n,p)$. \\
\rowcolor{tealpale} « Au moins un » succès & Passer par le complémentaire : $1-P(X=0)$. \\
\end{tabular}
\end{center}

\subsection{Les pièges qui coûtent des points}

\begin{remarque}[Erreurs relevées chaque année par les correcteurs]
\begin{itemize}[leftmargin=6mm, itemsep=0.5mm]
  \item Confondre $P(A\cap B)$ (probabilité de l'intersection) et $P_A(B)$ (probabilité conditionnelle) : $P(A\cap B)=P(A)\times P_A(B)$, ce ne sont pas la même chose.
  \item Utiliser la formule d'équiprobabilité alors que les issues \textbf{ne sont pas} équiprobables (vérifier l'énoncé).
  \item Oublier de vérifier que la somme des probabilités d'une loi vaut $1$.
  \item Confondre indépendance et incompatibilité : deux événements incompatibles ($A\cap B=\varnothing$) ne sont en général \textbf{pas} indépendants (sauf cas particulier).
  \item Dans la loi binomiale, oublier le coefficient binomial $\dbinom nk$ (compter le nombre d'arrangements des succès parmi les épreuves).
  \item Confondre $E(X)$ (valeur moyenne) et la valeur la plus probable de $X$.
\end{itemize}
\end{remarque}

% ------------------- RESUME -------------------
\vspace{4mm}
\begin{tcolorbox}[
  enhanced, boxrule=0pt, arc=2mm,
  interior style={left color=qtealdark, right color=qteal},
  left=6mm, right=6mm, top=4mm, bottom=4mm,
  fontupper=\color{white}\sffamily,
  title={\sffamily\bfseries\color{white}\ding{72}~~À retenir},
  colbacktitle=qtealdark, coltitle=white, boxrule=0pt
]
\begin{itemize}[leftmargin=6mm, label={\color{qamber}\ding{212}}, itemsep=1mm]
  \item Équiprobabilité : $P(A)=\frac{\text{card}(A)}{\text{card}(\Omega)}$ ; $P(\bar A)=1-P(A)$ ; $P(A\cup B)=P(A)+P(B)-P(A\cap B)$.
  \item Conditionnelle : $P_A(B)=\frac{P(A\cap B)}{P(A)}$ ; indépendance : $P(A\cap B)=P(A)P(B)$.
  \item Probabilités totales : $P(B)=\sum P(A_i)P_{A_i}(B)$ ; Bayes : $P_B(A_i)=\frac{P(A_i)P_{A_i}(B)}{P(B)}$.
  \item $E(X)=\sum x_iP(X=x_i)$ ; $V(X)=E(X^2)-E(X)^2$ ; $\sigma(X)=\sqrt{V(X)}$.
  \item Fonction de répartition : $F_X(t)=P(X\leq t)$, en escalier, croissante, saut $P(X=x_i)$ en $x_i$ ; $P(a<X\leq b)=F_X(b)-F_X(a)$.
  \item Loi binomiale $\mathcal{B}(n,p)$ : $P(X=k)=\binom nkp^k(1-p)^{n-k}$, $E(X)=np$, $V(X)=np(1-p)$.
\end{itemize}
\end{tcolorbox}

\end{document}
